% !TeX root = ../papers/nc-p3.tex
\section{対称成分の四つの境界因子}
\label{sec:symmetric-boundary-divisors}
$Z_{\mathrm{sym}}=\overline S$で、添字$e$以外の三因子を同時対称化できる
有限Frobenius候補のうち$S$に属さないものの閉包を$D_e$とする。
添字は固定し、四つの位置を同一視しない。

\begin{theorem}[三因子対称境界の余次元と局所被覆]\label{thm:four-boundary-divisors}
$D_0,D_1,D_2,D_3$は$Z_{\mathrm{sym}}\setminus S$の相異なる68次元既約成分である。
各因子の一般点は有限Frobeniusで非Koszulである。
各因子のある稠密開部分では、$\mathfrak X_6$は滑らかであり、その点の
$\mathfrak X_6$内の開近傍にある標準AS対象は全て$S$に属する。
この開部分には定理\ref{thm:nonkoszul-symmetric-boundary}のBoolean点の対応する巡回回転が入る。
\end{theorem}
\begin{proof}
$e=2$とし、$X=V^*$を対称な三因子の共通双対空間とする。
$\Gr(4,\Sym^2X)$の基点を持つ部分を$\Delta$とする。
基点$[p]$を指定したincidenceは$\mathbb P^3$上の$\Gr(4,9)$束で、既約かつ次元23である。
その射影像が$\Delta$である。Boolean点の基点はただ一つなので一般ファイバーは零次元であり、
$\Delta$は24次元Grassmann多様体の既約因子となる。

$B=\Sym(X)/(L)$で$\dim B_3=4$となる開部分上、$B_3^*$の射影フレームを取り、
対称な位置1と3に独立な$\PGL_4$変換を施す。位置0の座標は固定する。
このパラメータ空間は基底上の相対次元45の滑らかな束である。
一般の$L$はCIなので、一般の像は$S$に入る。従って全ての像は$Z_{\mathrm{sym}}$にある。
$\Delta$上でも$\dim B_3=4$となる開部分はBoolean点を含み、既約である。
この上の有限Frobenius条件は非空開である。実際、縮約の二対二階数は$6$以下であり
（三因子のうち二因子の積が$B_2$を経由する）、残る次元等号は
定理\ref{thm:frobenius-tensor-reconstruction}の開条件になる。
逆に三因子を対称化した有限Frobenius候補では、その連続三因子の交叉が$B_3^*$で
次元4である。系\ref{cor:partial-symmetric-exhaustive-split}から、$S$の外なら
$L\in\Delta$であり、こうした候補の全てがこの像に入る。

一般点では$S_1,S_3$は各々一次元で可逆行列を含む。
正規化した表示では両方が恒等行列を含み、Boolean点で次元1を達成するからである。
この二直線から位置1と3の射影基底変更を一意に戻せる。
三因子を同一視した後で、最初の二因子の縮約の消去子から$L$、残る因子の三次縮約から
$B_3^*$の射影フレームを回収できる。いずれも定数階数の核・像を取る正則な操作である。
従って$D_2$は次元$23+45=68$の既約閉部分である。
$S$は69次元既約な$Z_{\mathrm{sym}}$の開部分なので、$D_2$は境界の既約成分となる。
巡回回転で他の三位置についても同じ結論を得る。
異なる二つの像の一般部分が重なると、共通の因子を基準に取ったとき残る全因子との
交換に可逆な正証明書があり、完全対称化ができてしまう。これは基点枝の二分岐に反する。
同じ既約閉包を持つ稠密な構成可能集合は共通の開部分を持つため、四つの閉包は異なる。
一般部分の非Koszul性は定理\ref{thm:partial-symmetric-basepoint-exclusion}による。

最後にBoolean点で三分割の行列の左右核に対する接条件を全て課すと、
300行256列の有理行列の核次元は70である
（\path{computations/verify_boolean_boundary_divisor.py}）。
$\mathfrak X_6$は69次元成分$Z_{\mathrm{sym}}$を含み、射影接次元も69なので、
Boolean点で滑らか、通過する成分は$Z_{\mathrm{sym}}$だけである。
従ってこの点の近傍で$\mathfrak X_6$を$Z_{\mathrm{sym}}$に制限できる。
$Z_{\mathrm{sym}}$上では対称チャートからの特殊化により全ての正symmetrizer空間が非零である。
Boolean点で$S_1,S_3$は各々可逆な生成元を持つ一次元空間だから、
その性質と有限Frobenius性はこの点の開近傍で続く。
その近傍のAS点は二つの可逆証明書により全て対称CIとなる。
$D_2$との交叉は非空開、従って稠密である。
\end{proof}

\begin{proposition}[境界の横断座標と行列式の重複度]\label{prop:boolean-transverse-order}
Boolean点において$Z_{\mathrm{sym}}\setminus S$の台は滑らかな因子である。
位置2を三対称因子の一つへ同一視する正symmetrizer直線の局所生成行列$H$について、
$\det H$はこの因子に沿ってちょうど三重に消える。
一般点での四つの基本因子の付値は$(0,1,1,1)$である。
\end{proposition}
\begin{proof}
前証明の逆写像により、Boolean点の近傍は$L$、射影三次フレーム、二つの射影基底変更で
局所的に表示される。その中で$S$の補集合は$\Delta$の引戻しである。
Boolean点では最初の三二次式は八つの相異なる点を切り、第四式が零となる点は
$p=(0,0,0,1)$だけである。三式の勾配は$p$のアフィンチャートで独立で、
形式近傍では三式から$p$を一意に持ち上げられる。他の七点では第四式が単元である。
従って残る局所方程式は持ち上げた$p$での第四式の値一つであり、
第四式へ$t e_3^2$を加える曲線はこれを一次で横断する。

$S_2$もBoolean点で一次元なので局所的に直線をなす。
因子の一般点は単一被約基点$p'$を持ち、系\ref{cor:basepoint-rank-one-symmetrizer}から
その直線は階数1となる。従って単元成分で掃き出した後の$3\times3$ブロックの各成分は
滑らかな因子の局所方程式で割り切れ、行列式は少なくとも三重に消える。
正確な次数を上の曲線で計算する。八点を$v=(a,b,c,1)$、$s_v=a+b+c$と書く。
四次traceの正規化として
\[
 F_t=\sum_{v}\frac{(-1)^{s_v}t}{t+s_v}\,v^{\otimes4}
\]
を取れる（$s_v=0$の係数は1）。最初の三二次式は全八点で消え、第四式で縮約した値は
次数2以下の任意の単項式に対する符号付き立方体和で零となる。
$t=0$で$F_0=p^{\otimes4}$である。
三次フレームの正規化座標$012,133,233,333$を$m_j$とすると、
命題\ref{prop:partial-symmetrizer-quartics}の生成行列は$H_{ij}=F_t(i,m_j)$である。
八項の和の行列式を展開すると
\[
 \det H(t)=\frac{-2t^3}{(t+1)^3(t+2)^3(t+3)} .
\]
従って因子の局所方程式の三乗を除いた因子はBoolean点で単元である。
一般点の離散付値環では階数1への落ちと付値の和3から基本因子の主張を得る。
逆向きの射影同一視は逆行列のスカラー倍なので、一般点での極限階数は3となる。
計算コードは四次traceが三次フレームで表示されること、Boolean点の全12個の向き付き
symmetrizer空間が一次元で、通常の三位置相互は階数4、位置2からは1、位置2へは3であることも検算する。
\end{proof}

この四因子は三因子対称な有限Frobenius非正則境界の閉包を尽くす。
他の境界因子の有無や、これらの因子の全特殊点でのAS性は未決定である。
